% ============================================================================
% CONCLUSION GÉNÉRALE ET PERSPECTIVES
% ============================================================================

% ────────────────────────────────────────────────────────────────
\section*{Bilan des contributions}
\addcontentsline{toc}{section}{Bilan des contributions}

La première contribution majeure de ce projet est l'agent planificateur IA Guidni~: un agent
conversationnel basé sur LangGraph avec des outils appelables par le LLM, un pipeline RAG
multilingue utilisant les embeddings denses BAAI/bge-m3 (FR/EN/AR), une abstraction de
fournisseur LLM supportant Ollama, Groq et Google Gemini, une boucle de validation garantissant
la correction structurelle des plans, et un système de mémoire supportant les conversations
longues par résumé automatique à 8~messages. Le cycle de raisonnement LangGraph en quatre
nœuds (THINK → EXECUTE\_TOOL → VALIDATE → RESPOND) constitue une implémentation concrète du
paradigme ReAct appliqué à la planification de voyage. Ce travail constitue le premier agent
de planification de voyage conversationnel spécifiquement conçu pour le marché touristique de
Djerba avec un support multilingue complet français, anglais et arabe.

La deuxième contribution est le moteur de recommandation Hybrid Thompson Sampling~: un bandit
multi-bras bayésien exploitant la conjugaison Beta-Binomiale, une double segmentation sémantique
(clé globale de destination et clé utilisateur personnalisée), un mélange hybride pondéré par
un poids de personnalisation dynamique ($w = \min(0.8, N/40)$) résolvant élégamment le démarrage à froid,
un barème de calibrage comportemental des récompenses incluant un bonus financier de réservation, et un pipeline
de classement post-scoring en 10~étapes. L'efficacité du cache de 30~secondes en mémoire vive garantit des
lectures sous la milliseconde sans engorger la base de données. Combinés, ces éléments constituent un moteur
de recommandation hautement scalable, interprétable et adapté aux contraintes d'une place de marché
touristique~: apprentissage en ligne immédiat, exploration bayésienne naturelle et temps de calcul minimaux en $O(1)$.

La troisième contribution est l'intégration système~: les deux sous-systèmes partagent une
base de données PostgreSQL sans conflit grâce à une partition claire entre tables en lecture
seule (les ~tables miroirs du frontend Prisma/Next.js) et tables gérées par l'IA
(planificateur +  recommandation). La réindexation RAG incrémentale via des triggers SQL 
maintient l'index vectoriel toujours à jour sans intervention manuelle. Le tracker JavaScript relie le comportement client aux
mises à jour matricielles côté serveur de manière transparente, avec un mécanisme beacon
garantissant la capture des événements même à la fermeture de page. APScheduler exécute
tous les jobs de maintenance avec une planification timezone-aware (\texttt{Africa/Tunis})
et une protection contre les exécutions chevauchantes (\texttt{max\_instances=1}).

La quatrième contribution est la qualité d'ingénierie~: la base de code démontre plusieurs
pratiques d'ingénierie avancées. L'abstraction des fournisseurs LLM via une factory
uniforme (\texttt{get\_llm()}) permet le remplacement transparent du modèle sous-jacent.
La détection de changements par hachage MD5 dans le pipeline RAG évite les
recalculs d'embeddings inutiles, réduisant le coût de réindexation de 60\% dans les
scénarios de mise à jour par lots. Le cache matriciel en RAM avec verrou threading garantit
à la fois la performance (scoring $< 1$\,ms) et la sécurité thread. La limitation de
débit sur le traitement d'événements prévient la corruption des matrices par des scripts
automatisés. La journalisation structurée à plusieurs niveaux de verbosité facilite le
débogage et le monitoring en production.


% ────────────────────────────────────────────────────────────────
\section*{Perspectives et travaux futurs}
\addcontentsline{toc}{section}{Perspectives et travaux futurs}

La première perspective concerne l'affinage d'un LLM local sur les données du tourisme
tunisien. L'historique de conversations stocké dans la table \texttt{PlannerMessage}
constitue un corpus d'entraînement naturel~: paires (requête utilisateur, plan généré validé)
avec les métadonnées associées (outils utilisés, nombre d'itérations, score de validation).
Un fine-tuning supervisé ou par renforcement (RLHF) sur ce corpus produirait un modèle
domaine-spécifique comprenant la géographie locale (Houmt Souk, Midoun, Aghir, Er Riadh),
la culture tunisienne (Eid, été, hivernage), et les niveaux de prix réalistes de Djerba,
sans nécessiter d'appels API externes. Ce modèle local réduirait à la fois la latence
(pas de round-trip réseau) et le coût opérationnel (pas de facture API par requête).

La deuxième perspective est l'introduction d'un bandit contextuel bayésien multivarié ou de bandits neuronaux.
Bien que notre Thompson Sampling Hybride Beta-Binomial soit très robuste et performant, il n'intègre pas
directement de signaux continus dynamiques comme les données météorologiques instantanées ou l'heure de navigation.
Des architectures plus avancées (comme les bandits contextuels linéaires basés sur le Thompson Sampling gaussien
ou les Deep Contextual Bandits) modéliseraient ces effets fins par régression linéaire bayésienne ou réseaux de
neurones profonds. Bien que plus coûteuses en ressources de calcul et en volume de données, ces architectures
pourraient enrichir la personnalisation fine à mesure que le volume d'interactions de la plateforme grandit.

La troisième perspective est la construction d'une infrastructure d'expérimentation A/B.
L'architecture actuelle ne permet pas de mesurer l'impact causal de changements
algorithmiques sur les métriques métier réelles (taux de conversion, revenu par session,
durée de session). Un framework d'expérimentation assigne les utilisateurs à des groupes
traitement/contrôle, enregistre les métriques pour chaque groupe de manière indépendante,
et applique des tests statistiques (test-t, test séquentiel) pour détecter les
améliorations significatives. Ce framework permettrait de valider empiriquement les
hypothèses d'amélioration (nouvel algorithme, nouveau prompt, nouveaux tags) avant tout
déploiement en production.

La quatrième perspective est l'extension multi-régions. Le champ \texttt{segment}
dans la table \texttt{BanditArm} est nativement conçu pour supporter plusieurs zones
géographiques via la clé globale (par exemple, l'identifiant de destination)~; l'extension du système à Hammamet, Sousse, Tunis et d'autres destinations
tunisiennes ne nécessite que l'ajout de nouvelles données d'annonces et l'initialisation
des bras correspondants pour chaque nouveau segment géographique. Le même algorithme d'apprentissage s'applique à chaque destination. La
question ouverte est le choix entre des hyperparamètres $\alpha$ et $\beta$ initialisés de manière générique (partage de connaissances a priori
entre destinations) ou de manière totalement isolée par zone (démarrage à froid pour les nouvelles destinations). Une approche hybride avec
un a priori Beta basé sur la performance historique globale tunisienne est tout à fait envisageable.

La cinquième perspective est l'interface vocale et multimodale. L'intégration de la
reconnaissance vocale arabe permettrait aux utilisateurs tunisiens préférant l'expression
orale (notamment en darija tunisien) de planifier leur voyage sans saisie clavier. Côté
multimodal, le pipeline RAG pourrait accepter des photos de lieux ou de plats téléversées
par l'utilisateur comme requêtes visuelles, en utilisant des embeddings multimodaux (CLIP
ou des modèles vision-langage récents) pour retrouver des annonces visuellement similaires.
Ces extensions ouvriraient le système à une population d'utilisateurs plus large et
offriraient une expérience de planification encore plus naturelle.


% ────────────────────────────────────────────────────────────────
\section*{Mot de la fin}
\addcontentsline{toc}{section}{Mot de la fin}

Ce projet démontre que la combinaison de l'IA conversationnelle (agents LangGraph) et de
l'apprentissage en ligne (bandits contextuels) ouvre une nouvelle génération de plateformes
touristiques intelligentes~: simultanément personnalisées, adaptatives, et ancrées dans des
données réelles. L'agent planificateur Guidni montre qu'un LLM équipé d'outils bien conçus
et guidé par un prompt système rigoureux peut produire des plans de voyage cohérents,
multilingues et culturellement appropriés pour un marché spécifique. Le moteur Thompson Sampling Hybride montre
qu'un algorithme relativement simple --- basé sur des tirages de lois Beta probabilistes et un mélange
hybride élégant --- peut apprendre des préférences contextuelles et individuelles en temps réel, résolvant
le démarrage à froid et dépassant les baselines naïves avec seulement quelques centaines de sessions d'apprentissage. Le système Guidni est une
preuve de concept opérationnelle démontrant que cette combinaison est non seulement
techniquement réalisable, mais pratiquement déployable sur une infrastructure modeste~:
un serveur unique avec Ollama local, une base de données PostgreSQL standard, et un tracker
JavaScript sans aucune dépendance externe.
